\chapter{Experimental Scope and Reproducibility Notes}
\label{app:scope-notes}

\section{Measurement boundaries}

Physical-board evidence refers to records returned by a programmed Nexys 3 FPGA.
Post-route timing and resource values come from Xilinx ISE implementation reports.
RTL and software results are identified as simulations, and the OpenROAD results
refer to isolated physical-design kernels. These evidence types answer different
questions and are not combined into a single performance measure.

The reported kernel latency begins when a synchronous request is accepted and ends
when the FPGA asserts \texttt{done}. It excludes host processing, UART
serialization, JTAG programming, KMS acquisition, candidate generation, security
checks, and route installation.

\section{Repository milestone identifiers}

Some repository directories retain P0--P6 prefixes because the evidence was
archived in sequential engineering milestones. The directories identify the
software reference and contract foundation, literature verification, policy
training and freeze, controller verification, FPGA integration, and same-board
replay stages described by their accompanying completion records. These identifiers
are repository milestone labels and are not experimental variables. The scientific
variables used in this thesis are the hardware evaluation configurations H0--H4,
the research hypotheses RH1--RH5, scoring profiles $\Pi_0$--$\Pi_3$, and candidate
routes $\mathcal{R}_i$.

\section{Reproducibility records}

Every headline result can be traced to a version-controlled CSV, JSON summary,
implementation report, raw replay record, or physical-design report. Appendix
\ref{app:reproducibility} lists the relevant branches, commits, and paths. The
standalone file \texttt{FINAL\_THESIS\_NUMERICAL\_AUDIT.md} provides a compact
value-by-value index for the submission package.
